A single branch misprediction on a modern CPU can stall the execution pipeline for 15--20 cycles.
For a dataset of $10^7$ elements, a standard Quicksort evaluates roughly 240 million comparisons, each serving as a potential pipeline stall \cite{kaligosiHowBranchMispredictions2006}.
The motivation behind Adaptive Range Sorting (ARS) is to decouple the logical ordering of data from these costly arithmetic decisions.
By reframing sorting from a "decision tree" problem into a "signal projection" problem, ARS drastically minimizes the need for unpredictable branching.

\begin{figure}[ht]
\centering
\begin{tikzpicture}
\begin{axis}[
    ybar,
    width=0.98\linewidth, height=5.5cm,
    enlarge x limits=0.2,
    ylabel={Speedup vs. Quicksort (std) ($\times$)},
    symbolic x coords={Random i64, Gaussian i64, Duplicates, Strings},
    xtick=data,
    nodes near coords,
    nodes near coords align={vertical},
    ymin=0, ymax=4.5,
    grid=y,
    major grid style={dashed, gray!30},
    bar width=20pt,
    fill=teal!70,
    label style={font=\small},
    tick label style={font=\small}
]
\addplot coordinates {(Random i64, 2.39) (Gaussian i64, 1.78) (Duplicates, 1.28) (Strings, 3.36)};
\end{axis}
\end{tikzpicture}
\caption{The "Spatial Advantage": ARS Aero (Stable) Speedup across diverse data domains ($N=10^7$). The framework demonstrates significant gains in high-entropy and complex-type sorting while maintaining competitive parity on low-entropy duplicates.}
\label{fig:teaser}
\end{figure}

\subsection{Architectural Paradigm: Spatial Mapping}
ARS introduces a shift from comparative logic to \textbf{Spatial Mapping}.
By projecting arbitrary data types into a linear spatial context using a monotonic mapping function $\mathcal{K}$ (formally defined in Section \ref{sec:methodology}), the input sequence is treated as a discrete stochastic signal \cite{DrDobbsJournala}.
This allows the algorithm to analyze global data distributions and map values directly to hardware-aligned memory grids, achieving near-linear time performance across a broad class of stochastic distributions.
